% ============================================================================
\section{Conclusion}
\label{sec:conclusion}
% ============================================================================
This study presented an adaptive knowledge-based system that keeps payroll tax calculation current as government regulations and company policies change, including retroactively. Regulations and company policies form two layers of one executable knowledge graph whose rules record both their validity and the date they became known. A deterministic engine computes the tax. An LLM drafts the updates that deterministic gates and a behavioral consensus filter before an HC software engineer approves them. The system was replayed on the document streams of five Indonesian companies from 2023 to 2026 with an engineer approving its proposals. In four of the five companies it paid more employees correctly than a code-based update made 26 days after each issue. One incomplete installation of a labor agreement kept the pooled result below the pre-specified margin. A completeness gate added after the evaluation removed this failure in a replay of the recorded drafts. The engineer wrote updates by hand for 11 of 138 documents. A late upload of a retroactive rule left the earlier payroll reproducible in every scenario and produced the reference correction for the months already paid whenever the drafted update was correct.

The results also show where the approach needs care. Drafted updates of government regulations were mostly correct (35 of 39 prospective proposals), while drafted company policies often ended in an abstention or a silent error (24 of 30 prospective units). An update can also behave differently when it builds on the system's own earlier updates than when it is tested alone. Future work will test the completeness gate on new documents and repeat the chains with several runs and more companies. It will also generate review cases from the company's current knowledge base and from the months in which an update takes effect.
